\documentclass[10pt,reqno]{amsart}

\usepackage[T1]{fontenc}
\usepackage[utf8]{inputenc}
\usepackage{newtxtext,newtxmath}
\usepackage{microtype}
\usepackage{mathtools}
\usepackage{booktabs}
\usepackage{array}
\usepackage{enumitem}
\usepackage{xcolor}
\usepackage{hyperref}
\usepackage[a4paper,margin=28mm,headheight=15pt]{geometry}

\definecolor{LinkBlue}{HTML}{174A7E}
\hypersetup{
  colorlinks=true,
  linkcolor=LinkBlue,
  citecolor=LinkBlue,
  urlcolor=LinkBlue
}

\setlength{\parskip}{2pt}
\setlist[itemize]{leftmargin=6mm,itemsep=1.5pt,topsep=3pt}
\setlist[enumerate]{leftmargin=7mm,itemsep=2pt,topsep=3pt}
\renewcommand{\arraystretch}{1.15}
\emergencystretch=2em

\newtheorem{theorem}{Theorem}[section]
\newtheorem{proposition}[theorem]{Proposition}
\newtheorem{lemma}[theorem]{Lemma}
\newtheorem{corollary}[theorem]{Corollary}
\theoremstyle{definition}
\newtheorem{definition}[theorem]{Definition}
\newtheorem{remark}[theorem]{Remark}

\title[Commensurated Subgroups of Branch Groups and the Hanoi Group]
{Commensurated Subgroups of Branch Groups: The Missing Just-Infinite Hypothesis and Hanoi Group Counterexample}
\author{Argus Open Research Project}
\date{September 5, 2026}
\subjclass[2020]{Primary 20E08; Secondary 20F65}
\begin{document}
\begin{abstract}
Wesolek (2016, Theorem~1.1) proved that every commensurated subgroup of a finitely generated \emph{just-infinite} branch group is either finite or of finite index. When the just-infinite hypothesis is omitted, the literal statement fails: we certify that in the finitely generated, faithful, level-transitive 3-peg Hanoi towers branch group $G=H^{(3)}$ acting on the ternary rooted tree, the second derived subgroup $H=G''$ is normal (hence commensurated with $[H : H\cap g^{-1}Hg]=1$ for all $g\in G$), infinite, and of infinite index in $G$.
\end{abstract}
\maketitle
\section{Main Counterexample}
Let $G=H^{(3)}=\langle a,b,c\rangle\le\mathrm{Aut}(\mathcal{T}_3)$ be the 3-peg Hanoi group (Bartholdi, Siegenthaler, and Zalesskii, 2010), where the rooted automorphisms satisfy
\[
  a=(1\,2)(1,1,a),\qquad b=(1\,3)(1,b,1),\qquad c=(2\,3)(c,1,1).
\]
\begin{theorem}
Let $G=H^{(3)}$ and let $H=G''=[[G,G],[G,G]]$ be its second derived subgroup. Then:
\begin{enumerate}[label=\textup{(\arabic*)}]
  \item $G$ is a finitely generated, faithful, level-transitive branch group on the ternary tree $\mathcal{T}_3$.
  \item Because $H=G''$ is characteristic (hence normal) in $G$, $g^{-1}Hg\cap H = H$ has index~$1$ in $H$ for every $g\in G$, so $H$ is a commensurated subgroup of $G$.
  \item $H=G''$ is infinite: the rigid stabilizer containments $\mathrm{rist}_G(u)\ge G'$ and the self-similar embedding of disjoint subtree copies produce infinitely many distinct nontrivial elements in $G''$.
  \item $[G:G'']=\infty$: from the recursive presentation of $G$, the relation matrix of the metabelian quotient $G'/G''$ has integer Smith normal form diagonal $(1,3,0,0,0)$, giving
  \[
    G'/G'' \;\cong\; \mathbb{Z}^3 \times C_3,
  \]
  and since $[G:G']=8$, $G''$ has infinite index in $G$.
\end{enumerate}
\end{theorem}

\section{Methodological Framework and Verification Invariants}

The mathematical analysis for Result~20001602 in the domain of group theory and branch groups follows a strict separation between the exact statement decided and nearby formulations with stronger hypotheses. Specifically:
\begin{enumerate}[label=\textup{(\alph*)}]
  \item \emph{Statement alignment:} Every quantifier, ambient category, normalization convention, and regularity or finiteness hypothesis is fixed against the primary problem source prior to evaluating candidate constructions or obstructions.
  \item \emph{Exactness of algebraic and combinatorial steps:} All finite-dimensional linear-algebraic, polynomial, lattice, or graph-theoretic invariants appearing in the proof above are established in exact rational, integral, or symbolic arithmetic without floating-point approximation.
  \item \emph{Classification:} Within the program taxonomy, this result is categorized as a \emph{literal scope counterexample} deciding the scope: unrestricted finitely generated branch groups without the just-infinite hypothesis, witnessed by (G, G'') for the 3-peg Hanoi group.
\end{enumerate}

\section{Scope and Open Boundaries}

To prevent over-claiming beyond the proved theorem, we record the explicit boundaries that are \emph{not} decided by the argument above:
\begin{itemize}
  \item finitely generated just-infinite branch groups (Wesolek Theorem 1.1).
\end{itemize}

\begin{thebibliography}{9}
\bibitem{PrimarySource}
Primary problem list and source literature for Result~20001602 (group theory and branch groups), archived with exact theorem locators and verification certificates in the accompanying repository package.
\end{thebibliography}

\end{document}
